\chapter{Procesor potokowy}
\label{ch:pipe}

\metryczka{przerobić procesor jednocyklowy na klasyczny 5-etapowy potok z obsługą
wszystkich hazardów i zmierzyć, co to dało. Najtrudniejszy i najciekawszy
rozdział kursu.}{4--6 wieczorów}{08-pipeline/}

% Diagram potoku: \instr{wiersz}{pierwszy takt}{etap1,etap2,...}{mnemonik}
\newcommand{\instr}[4]{%
  \node[anchor=east, font=\ttfamily\scriptsize] at (0.9,-#1) {#4};
  \foreach \s [count=\k from 0] in {#3} {
    \ifx\s\empty\else
    \pgfmathsetmacro{\xx}{#2+\k}
    \node[draw, thick, minimum width=9.5mm, minimum height=5mm, inner sep=0pt,
      font=\sffamily\scriptsize, fill=akcentjasny] at (\xx,-#1) {\s};
    \fi
  }}
\newcommand{\banka}[2]{% bańka: wiersz, takt
  \node[draw, thick, dashed, minimum width=9.5mm, minimum height=5mm, inner sep=0pt,
    font=\sffamily\tiny, fill=white, text=szary] at (#2,-#1) {bańka};}
\newcommand{\takty}[1]{%
  \foreach \t in {1,...,#1} { \node[font=\sffamily\scriptsize, text=szary] at (\t+1,0.6) {\t}; }}

\section{Po co potok?}

Czas wykonania programu („żelazne prawo” wydajności procesora):
\[ \text{czas} \;=\; \text{liczba instrukcji} \times \text{CPI} \times T_{clk}. \]

Procesor jednocyklowy ma CPI = 1, ale bardzo długi takt. Potok tnie drogę
instrukcji rejestrami na 5 krótszych etapów. Każda instrukcja nadal potrzebuje 5
etapów, więc \emph{opóźnienie} pojedynczej instrukcji się nie zmienia. Ale \textbf{w
każdym takcie kończy się jedna instrukcja}, a takt może być kilka razy krótszy.

\begin{figure}[h]
\centering
\begin{tikzpicture}[x=10.5mm, y=7mm]
  \takty{8}
  \instr{1}{2}{IF,ID,EX,MEM,WB}{add x1,..}
  \instr{2}{3}{IF,ID,EX,MEM,WB}{sub x2,..}
  \instr{3}{4}{IF,ID,EX,MEM,WB}{lw x3,..}
  \instr{4}{5}{IF,ID,EX,MEM,WB}{beq ..}
  \draw[pulapka, thick, rounded corners] (5.5,0.3) rectangle (6.5,-4.4);
  \node[etykieta, text=pulapka, align=center] at (6,-4.9) {w takcie 5 pracują\\wszystkie etapy};
\end{tikzpicture}
\caption{Idealny potok: od piątego taktu kończy się jedna instrukcja na takt.}
\label{fig:pipe}
\end{figure}

Cena: instrukcje „zachodzą” na siebie i zależności między nimi trzeba obsłużyć
w sprzęcie. To są \textbf{hazardy}.

\section{Etapy i rejestry potoku}

\begin{center}
\small
\begin{tabularx}{\linewidth}{@{}l X X@{}}
\toprule
\textbf{Etap} & \textbf{Co robi} & \textbf{Rejestr na końcu etapu} \\
\midrule
IF & \code{imem\_addr = pc}, pobranie instrukcji, \code{pc += 4} & IF/ID: \code{d\_valid, d\_pc, d\_instr} \\
ID & dekoder, immediate, odczyt rejestrów, wykrycie load-use & ID/EX: sterowanie, \code{rs1\_val, rs2\_val, imm}, numery rs1/rs2/rd \\
EX & forwarding, ALU, \code{branch\_cmp}, cel skoku, \textbf{redirect} & EX/MEM: wynik, dana do store, rd \\
MEM & LSU i pamięć danych & MEM/WB: wartość do zapisu, rd \\
WB & zapis do banku rejestrów, \textbf{commit} & --- \\
\bottomrule
\end{tabularx}
\end{center}

\begin{figure}[h]
\centering
\begin{tikzpicture}[x=1cm, y=1cm, every node/.append style={font=\sffamily\scriptsize}]
  \foreach \x/\n in {0/IF, 2.9/ID, 5.8/EX, 8.7/MEM, 11.6/WB} {
    \node[blok, minimum width=17mm, minimum height=14mm] (\n) at (\x,0) {\n};
  }
  \foreach \x/\n/\r in {1.45/{IF/ID}/ra, 4.35/{ID/EX}/rb, 7.25/{EX/MEM}/rc, 10.15/{MEM/WB}/rd} {
    \node[draw, thick, fill=zadaniejasny, minimum width=3.5mm, minimum height=22mm] (\r) at (\x,0) {};
    \node[etykieta] at (\x,1.4) {\n};
  }
  \draw[->, thick] (IF) -- (ra); \draw[->, thick] (ra) -- (ID);
  \draw[->, thick] (ID) -- (rb); \draw[->, thick] (rb) -- (EX);
  \draw[->, thick] (EX) -- (rc); \draw[->, thick] (rc) -- (MEM);
  \draw[->, thick] (MEM) -- (rd); \draw[->, thick] (rd) -- (WB);
  % redirect
  \draw[->, thick, pulapka] (EX.north) -- ++(0,1.25) -| node[above, pos=0.25]{redirect (skok wzięty)} (IF.north);
  % forwarding
  \draw[->, thick, zadanie] ($(rc.south)$) -- ++(0,-0.6) -| node[below, pos=0.3]{EX/MEM → EX} ($(EX.south)+(0.3,0)$);
  \draw[->, thick, zadanie] ($(rd.south)$) -- ++(0,-1.1) -| node[below, pos=0.2]{MEM/WB → EX} ($(EX.south)+(-0.3,0)$);
  % zapis rejestru
  \draw[->, thick, akcent] (WB.north) -- ++(0,1.85) -| node[above, pos=0.25]{zapis rd (+ bypass WB → ID)} (ID.north);
\end{tikzpicture}
\caption{Pięć etapów, cztery rejestry potoku. Czerwony: redirect przy skoku.
Zielone: ścieżki forwardingu. Niebieski: zapis do banku rejestrów.}
\label{fig:pipeline}
\end{figure}

\paragraph{Konwencje} (zgodne ze szkieletem \plik{rtl/core.sv}): przedrostek
\code{d\_} oznacza sygnał w ID, \code{e\_} w EX, \code{m\_} w MEM, a \code{w\_} w WB.
Każdy etap ma bit \code{*\_valid}. \textbf{Bańka} (\emph{bubble}) to etap z
\code{valid = 0} i wyzerowanymi sygnałami, które mają skutki uboczne
(\code{reg\_write}, \code{mem\_write}, \code{branch}, \code{jump} \dots). Przepływa
przez potok jak \code{nop}, ale nie jest raportowana jako commit.

\paragraph{Commit z WB.} \code{commit\_valid = w\_valid}, a opis instrukcji
(\code{pc}, \code{instr}) jedzie przez cały potok. Dzięki temu ślad jest identyczny
jak w emulatorze i \code{make cosim} działa bez zmian.

\section{Hazardy danych}

\begin{figure}[h]
\centering
\begin{tikzpicture}[x=10.5mm, y=7mm]
  \takty{6}
  \instr{1}{2}{IF,ID,EX,MEM,WB}{add x1,x2,x3}
  \instr{2}{3}{IF,ID,EX,MEM,WB}{sub x4,x1,x5}
  \draw[->, very thick, pulapka] (6,-1.3) .. controls (5.5,-1.7) and (4.6,-1.6) .. (4.1,-1.75);
  \node[etykieta, text=pulapka, anchor=west] at (7.7,-1) {x1 zapisany na końcu WB (takt 5)};
  \node[etykieta, text=pulapka, anchor=west] at (7.7,-2) {a czytany w ID już w takcie 3!};
\end{tikzpicture}
\caption{Hazard RAW (\emph{read after write}).}
\end{figure}

Rozwiązania, od najprostszego:
\begin{enumerate}
  \item \textbf{Programowo}: \code{nop}-y między zależnymi instrukcjami. Tak
    działa \code{make progs PAD=n} (\code{rvtool --pad-nops}) na etapie 8.1.
  \item \textbf{Bypass WB → ID}: gdy WB w tym samym takcie zapisuje rejestr, który
    ID czyta, weź wartość prosto z WB (bank rejestrów oddałby starą).
  \item \textbf{Forwarding do EX}: wynik \code{add} istnieje już na końcu EX.
    Następna instrukcja potrzebuje go na \emph{początku} swojego EX, czyli takt
    później. Przekaż go z rejestru EX/MEM albo MEM/WB prosto na wejście ALU.
  \item \textbf{Stall} przy \emph{load-use} (niżej).
\end{enumerate}

\begin{figure}[h]
\centering
\begin{tikzpicture}[x=10.5mm, y=7mm]
  \takty{6}
  \instr{1}{2}{IF,ID,EX,MEM,WB}{add x1,x2,x3}
  \instr{2}{3}{IF,ID,EX,MEM,WB}{sub x4,x1,x5}
  \instr{3}{4}{IF,ID,EX,MEM,WB}{or  x6,x1,x7}
  \draw[->, very thick, zadanie] (4.45,-1.2) -- (4.55,-1.8);
  \draw[->, very thick, zadanie] (5.45,-1.2) .. controls (5.6,-2.2) .. (5.55,-2.8);
  \node[etykieta, text=zadanie, anchor=west] at (8.8,-1.6) {EX/MEM → EX (odległość 1)};
  \node[etykieta, text=zadanie, anchor=west] at (8.8,-2.6) {MEM/WB → EX (odległość 2)};
\end{tikzpicture}
\caption{Forwarding: wynik z rejestru potoku wraca prosto na wejście ALU.}
\end{figure}

\begin{sv}
always_comb begin
  if (m_valid && m_reg_write && m_rd != 0 && m_rd == e_rs1)  fwd_a = m_result; // najnowszy!
  else if (w_valid && w_reg_write && w_rd != 0 && w_rd == e_rs1) fwd_a = w_val;
  else                                                     fwd_a = e_rs1_val;
  // to samo dla fwd_b; fwd_b jest też daną dla store!
end
\end{sv}

\begin{pulapka}[Trzy szczegóły forwardingu]
(1) \textbf{Kolejność}: gdy oba etapy piszą ten sam rejestr, wygrywa
\emph{nowszy}, czyli EX/MEM (test \code{11\_hazards}, podtest 4).
(2) \textbf{\code{rd != 0}}: zapis do \code{x0} nie może być forwardowany
(podtest 8). (3) Dana do \textbf{store} (\code{rs2}) też musi przejść przez
forwarding.
\end{pulapka}

\subsection*{Load-use: jedyny przypadek, gdy trzeba czekać}

\begin{figure}[h]
\centering
\begin{tikzpicture}[x=10.5mm, y=7mm]
  \takty{7}
  \instr{1}{2}{IF,ID,EX,MEM,WB}{lw  x1,0(x2)}
  \instr{2}{3}{IF,ID,ID,EX,MEM,WB}{add x3,x1,x1}
  \banka{3}{5}
  \node[etykieta, anchor=west] at (5.6,-3) {bańka wstawiona do ID/EX};
  \draw[->, very thick, zadanie] (5.45,-1.2) -- (5.55,-1.8);
  \node[etykieta, text=zadanie, anchor=west] at (8.1,-1.4) {dana z pamięci jest dopiero po MEM};
\end{tikzpicture}
\caption{Load-use: instrukcja zależna czeka jeden takt w ID, a potem dostaje daną przez forwarding z MEM/WB.}
\end{figure}

Wykrywasz to w ID: w EX jest load, a jego \code{rd} to \code{rs1} albo \code{rs2}
instrukcji w ID. Wtedy \textbf{PC i IF/ID stoją}, a do ID/EX idzie bańka.

\section{Hazardy sterowania}

Skok rozstrzyga się w EX. Do tego czasu IF i ID zdążyły pobrać dwie instrukcje
„za skokiem”. Jeśli skok jest wzięty: \code{pc = cel}, a instrukcje w IF/ID i
ID/EX trzeba \textbf{unieważnić} (\emph{flush}).

\begin{figure}[h]
\centering
\begin{tikzpicture}[x=10.5mm, y=7mm]
  \takty{7}
  \instr{1}{2}{IF,ID,EX,MEM,WB}{beq (wzięty)}
  \instr{2}{3}{IF,ID}{inst+4}
  \instr{3}{4}{IF}{inst+8}
  \instr{4}{5}{IF,ID,EX,MEM,WB}{cel skoku}
  \draw[pulapka, very thick] (2.55,-1.75) -- (3.45,-2.25) (2.55,-2.25) -- (3.45,-1.75);
  \draw[pulapka, very thick] (3.55,-1.75) -- (4.45,-2.25) (3.55,-2.25) -- (4.45,-1.75);
  \draw[pulapka, very thick] (3.55,-2.75) -- (4.45,-3.25) (3.55,-3.25) -- (4.45,-2.75);
  \node[etykieta, text=pulapka, anchor=west] at (6.8,-2.5) {unieważnione: kara 2 takty};
\end{tikzpicture}
\caption{Wzięty skok: dwie instrukcje pobrane „na ślepo” zamieniają się w bańki.}
\end{figure}

Strategia „zakładaj, że skok nie będzie wzięty” to najprostsza predykcja skoków.
\textbf{Priorytety}: gdy w jednym takcie jest i redirect, i stall, wygrywa redirect,
bo czekająca instrukcja i tak zostanie unieważniona.

\section{Plan pracy}

Szkielet \plik{rtl/core.sv} ma sekcje dla etapów. Pracuj etapami i commituj po
każdym, który działa.

\begin{center}
\small
\begin{tabularx}{\linewidth}{@{}l X l@{}}
\toprule
\textbf{Etap} & \textbf{Co dodajesz} & \textbf{Test} \\
\midrule
8.1 & rejestry potoku, bity \code{valid}, redirect i flush przy skokach & \code{make progs PAD=4} \\
8.2 & forwarding EX/MEM→EX i MEM/WB→EX (oba operandy i dana do store), bypass WB→ID & \code{make progs PAD=1} \\
8.3 & stall load-use & \code{make progs} \\
8.4 & weryfikacja & \code{make cosim}, \code{make fuzz N=200} \\
8.5 & pomiary: CPI, synteza & \code{make synth T=core} \\
\bottomrule
\end{tabularx}
\end{center}

\code{PAD=n} wstawia $n$ instrukcji \code{nop} \emph{przed} każdą instrukcją. Etykiety
wskazują na pierwszy z nich, żeby adresy powrotu \code{jal} się zgadzały. Przy
\code{PAD=4} żadna zależność danych nie jest groźna.

\section{Zadania}

\begin{zadanie}{8.1 --- Potok bez obsługi hazardów danych}
\code{make progs PAD=4}. Jaki jest najmniejszy PAD, przy którym Twój potok z
etapu 8.1 przechodzi wszystkie testy? Wylicz to z diagramu potoku, zanim
sprawdzisz. Jak zmienia się ta liczba po dodaniu samego bypassu WB→ID?
\end{zadanie}

\begin{zadanie}{8.2 --- Forwarding i bypass}
\code{make progs PAD=1}. Dlaczego przy \code{PAD=0} nadal coś pada? Które testy i
dlaczego akurat te?
\end{zadanie}

\begin{zadanie}{8.3 --- Stall load-use}
\code{make progs}: 12 × PASS bez żadnych \code{nop}-ów.
\end{zadanie}

\begin{zadanie}{8.4 --- Weryfikacja}
\code{make cosim}, \code{make fuzz N=200}, a także
\code{make -C ../06-weryfikacja fuzz CORE=08 N=1000 SIM=verilator} i
\code{make -C ../07-c-bare-metal test CORE=08}. Fuzzer jest tu naprawdę cenny:
losowe programy tworzą kombinacje zależności, skoków i loadów, których nie ma w
ręcznych testach.
\end{zadanie}

\begin{zadanie}{8.5 --- Pomiary}
Porównaj CPI z \code{make progs} z rozdziałem 5 (zawsze 1,000). Na rozwiązaniu
wzorcowym \code{09\_sort} ma CPI ok.\ 1,5, a \code{bench} z rozdziału 7 ok.\ 1,3. Ile
taktów zabierają skoki, a ile load-use? Potem \code{make synth T=core} tu i w
rozdziale 5. Który etap ma teraz najdłuższą ścieżkę? (Podpowiedź: EX, czyli
forwarding, ALU, decyzja o skoku i multiplekser PC.)
\end{zadanie}

\begin{zadanie}{8.6 --- Dla ambitnych}
Skoki rozstrzygane w ID (kara 1 takt, ale potrzebny forwarding do ID i nowy
stall); statyczna predykcja BTFN (skoki wstecz zakładamy jako wzięte);
predyktor dynamiczny z 2-bitowymi licznikami i BTB. Mierz CPI na \code{bench}.
\end{zadanie}

\begin{gotowe}
\code{make progs} → 12 × PASS przy \code{PAD=0}, \code{make cosim} → wszystko
ZGODNE, \code{make fuzz N=500} → 0 różnic, a \code{bench} z rozdziału 7 działa na
\code{CORE=08}.
\end{gotowe}

\begin{kontrolne}
  \item Dlaczego potok skraca czas wykonania programu, choć każda instrukcja trwa tyle samo albo dłużej?
  \item Narysuj diagram potoku dla \code{lw x1,0(x2); sw x1,4(x2)}. Czy potrzebny jest stall?
  \item Dlaczego przy forwardingu EX/MEM ma pierwszeństwo przed MEM/WB?
  \item Ile taktów traci potok na wziętym skoku i jak to zmniejszyć?
  \item Co się stanie, jeśli przy stallu zapomnisz zatrzymać rejestr IF/ID (tylko PC)?
\end{kontrolne}
